\item
  \textbf{nonlinear design effect}: tanh plus a quadratic term in the quality proxy;
\item
  \textbf{heteroscedastic noise}: the standard deviation depends on the site;
\item
  \textbf{design-by-cell-type interaction}: two feature blocks whose mixing weight depends on the site, with the biological effect confined to one block;
\item
  \textbf{20\% case rate, matched thresholds}: the collection boundary is moved to the same quantile as the diagnosis threshold;
\item
  \textbf{20\% case rate, offset boundary}: the collection boundary stays at the median, so the threshold and the boundary are offset;
\item
  \textbf{different acquisition in the target}: the external cohort is acquired under a different design mechanism. The two 20\% arms are counted separately because they behave differently, and separating them is what showed the earlier ``rare diagnosis blunts the checks'' result to be a threshold artefact (Section 3.3).
\end{enumerate}

Critically, the external target reuses the \textbf{source's} biological loading vector. Regenerating it makes every transfer chance-level by construction and tests nothing; only the design mechanism differs between source and target in the shift arm.

A supplementary arm quantifies what admitting a disease-caused covariate does. A severity variable was generated as 1.2·(y − ȳ) + noise and added to the design matrix. Everything else was held fixed.

\subsection{9.10. Decision-tree walkthrough}\label{decision-tree-walkthrough}

Three cohorts were traced through the tree with all evidence recomputed, and every branch value is released in Table S9. Collection strata used exactly the definition in Section 9.6, so that Q2 and Q4 refer to the same partition. Restriction was evaluated in the largest stratum containing at least five donors of each class; where no stratum qualifies, that is reported rather than worked around by coarsening the strata. Matched random subsets of identical size and case/control composition were drawn 20 times from the full cohort for comparison.

\subsection{9.11. Recovery of GSE135779 design metadata}\label{recovery-of-gse135779-design-metadata}

Batch, collection year, age, sex, race, ethnicity and clinical variables were read from the original Supplementary Table 1b, and per-library sequencing metrics from Supplementary Table 1c {[}1{]}. Study names were linked to analysis donor identifiers through the GEO title/accession map; all 44 donors mapped uniquely. The restoration script writes the donor table, batch-by-label, year-by-label and batch-by-year cross-tabulations, and SHA256 hashes of every source file. For GSE174188, per-cell \texttt{Processing\_Cohort} was read from the distributed CELLxGENE object and reduced to donor fractions over four waves. The dominant wave is the maximum-fraction wave; a pure-wave donor has at least 99\% of analysed cells in one wave.

\subsection{9.12. Representations}\label{representations}

\textbf{Frozen Geneformer.} Cell embeddings used Geneformer V2-316M at repository revision \texttt{04c2b2e84da7c0f385c3f9ad8f3ec24bab6650e5} {[}4{]}; checkpoint, configuration, token dictionary and gene-median dictionary hashes are in the methods manifest. Gene identifiers were version-trimmed and mapped to the V2 vocabulary. Within each cell the ranking value was \texttt{raw\ count\ /\ total\ count\ ×\ 10,000\ /\ Genecorpus-104M\ gene\ median}, and genes were sorted descending. Sequences used V2 start/end token IDs 2 and 3, padding ID 0 and maximum length 4096, leaving at most 4094 ranked genes. The model output was the final \texttt{last\_hidden\_state}; token position 0 was kept as the 1,152-dimensional cell vector. The encoder was frozen. Donor representations are coordinate-wise means of sampled cell vectors, sampled without replacement with integer seed 1, capped at 500 cells per donor for GSE135779 and 1,000 for GSE174188 and GSE285773. The embedding environment used Python 3.10, PyTorch 2.5.1 and Transformers 4.46.3.

\textbf{Pseudobulk.} We use pseudobulk rather than a learned latent space as the donor-level representation, so that the comparison between representations is not itself mediated by a model fitted with batch covariates {[}5{]}. Stored donor-by-gene values are \texttt{log1p(CPM)}. Inside every outer training fold, gene variance was computed on training donors, the 4,000 highest-variance genes were selected, and coordinates were standardised by training-donor mean and standard deviation. HVG pseudobulk uses these values directly. PCA pseudobulk fits up to 30 components on the training values and applies the fitted map to held-out donors. For cross-cohort transfer, trimmed gene identifiers were intersected before source-only HVG selection; feature means, variances, selected genes, scaler, PCA, classifier and regularisation were all fitted on source donors and applied unchanged to the target.

\subsection{9.13. Fold-contained residualisation and its variants}\label{fold-contained-residualisation-and-its-variants}

Unadjusted and residualised models share outer folds, feature construction, scaling, PCA and classifier selection. Inside each outer training fold, design variables were encoded as in Section 9.5 and a multivariate ridge regression (\texttt{alpha} 1, intercept fitted and \textbf{not} penalised) mapped design to every raw representation coordinate. Training and held-out coordinates were replaced by observed minus predicted values, after which the identical pipeline was fitted to training residuals and applied to held-out residuals. The diagnosis was never given to the residualiser.

Three further arms were run on GSE135779. A multivariate random-forest residualiser (96 trees, depth 4, minimum leaf 3) fitted only on outer-training donors, with pseudobulk genes pre-selected by training-donor variance. A batch-location arm subtracting training-batch mean offsets from training and held-out features; because it applies no empirical-Bayes scale shrinkage we describe it as ComBat-style location adjustment rather than ComBat. And an overlap-weighting arm fitting a training-only logistic diagnosis propensity model from the full design, clipping propensities to 0.025-0.975 and training the representation classifier with overlap weights.

For the central GSE174188 comparison, residualisation loss was paired by split seed. We bootstrapped the 20 paired repeat-level losses 5,000 times and subtracted the locked matched-restriction discrepancy. These intervals quantify split sensitivity conditional on the analysed donors; they are not population confidence intervals.

To probe confound leakage with a nonlinear downstream learner {[}28{]}, 100 Gaussian simulations used no biological effect, a technical effect of one, and design-diagnosis association 0.75. Random-forest prediction was compared on raw features, globally residualised features and fold-contained residualised features.

\subsection{9.14. Restriction and matched donor controls}\label{restriction-and-matched-donor-controls}

GSE174188 was restricted to dominant wave 4 and separately to pure wave 4. GSE135779 excluded the all-case batch B1 and retained B2 to B6, each containing both labels. The complete internal pipeline was re-run after restriction. For every observed stratum, 20 donor subsets were drawn without replacement from the full cohort with identical donor count and case/control count, each using one prespecified split seed. The observed-minus-matched difference asks whether the stratum is harder than expected from reduced donor number and label composition alone.

\subsection{9.15. Cell-type composition}\label{cell-type-composition}

The GSE174188 object was restricted to the same CD4-positive alpha-beta T-cell population and 261 donors as the molecular analysis. Per-donor counts were formed for \texttt{T4\_naive}, \texttt{T4\_em}, \texttt{T4\_reg} and all other or unassigned labels. We evaluated raw closed proportions and centred-log-ratio coordinates; for CLR, 0.5 was added to each donor-category count before closure, the donor mean log abundance was subtracted and the final coordinate omitted. Each representation was evaluated with and without log total CD4-cell yield, under the same repeated donor pipeline, with processing wave predicted one against the rest.

\subsection{9.16. Transfer and learned pooling}\label{transfer-and-learned-pooling}

Cross-batch analyses trained on waves 2 and 3 and evaluated wave 4, then reversed. Cross-cohort transfer fitted every statistic on the source cohort; target labels were accessed only after prediction. Target AUC intervals used 5,000 case/control-stratified donor bootstrap resamples and the percentile method. Paired AUC comparisons used DeLong tests on identical target donors {[}23{]} with Benjamini-Hochberg correction inside declared method families {[}24{]}.

DeepSets, gated-attention multiple-instance learning and pooling by multi-head attention operated on cap-500 frozen cell embeddings {[}20-22{]}. Cell standardisation, a whitened source PCA32 projection, network weights and hyperparameters were all source-fitted. Hidden widths were 16 or 32, weight decay 0.01 or 0.1, and Adam used learning rate 0.02 for 150 epochs; five-fold source-donor cross-validation selected the configuration, three initialisations were fitted on all source donors, and target probabilities were averaged. An ordinary donor mean in the identical PCA32 space isolates the contribution of learned pooling. The PCA32 projection used for hyperparameter selection was fitted once on all source cells before source-donor folds and never used target cells or labels. Source-internal pooling AUC is therefore not an unbiased generalisation estimate, and only independent-target and paired target-donor comparisons support the learned-pooling conclusion.

\subsection{9.17. Software and reproducibility}\label{software-and-reproducibility}

The donor-level analyses used Python 3.11-3.13 with NumPy 2.3-2.5, pandas 2.3-3.0, scikit-learn 1.9.0, PyArrow and joblib; exact per-run versions are recorded in each output manifest. Figures were produced in R 4.6.0 with ggplot2 4.0.3 and patchwork 1.3.2. The screen, the extended simulation, the permutation calibration and the decision-tree walkthrough were run on a 16-vCPU ARM container; all other analyses were run locally. Input hashes, seeds, parameter grids, donor-level predictions, complete null distributions and output manifests accompany the analysis.

Random seeds are derived with \texttt{hashlib.blake2b} from a recorded master seed and the cell identity. Python's built-in \texttt{hash} is not used: it is salted per process for strings, so it would not reproduce across interpreter sessions. The derived seed for every cell is written into the released result tables. No script contains a machine-specific path: the figure scripts locate the project root from their own file position or from \texttt{RHEUMLENS\_ROOT}, and the analysis scripts take input and output paths as arguments.

Everything downstream of the compute-heavy runs is rebuilt by one command, \texttt{bash\ tools/build\_all.sh}: the supplementary tables, Table 1, all fifteen figures, the typeset manuscript, and two checks that fail with a non-zero exit status. The first re-derives numbers quoted in the text from the result tables and requires each to appear verbatim. The second checks the length of the abstract and author summary, that the abstract is unstructured, that figures and tables are cited in the order they are declared, and that each declared item is cited at least once. It also refuses any placeholder, retired term or withdrawn phrase, and it looks inside the rendered figures and the released table headers, where a check on the manuscript text alone cannot see them. The submission package is assembled by a further script from the same outputs, so a figure number cannot differ between the manuscript and the files.

\subsection{9.18. Use of generative AI}\label{use-of-generative-ai}

Anthropic Claude (models Claude Opus 4.1 and Claude Opus 5), accessed through the Claude Code command-line interface, was used during this work in the following places and no others.

\begin{itemize}
\tightlist
\item
  \textbf{Code review and refactoring} of the analysis scripts in \texttt{scripts/cohorts/}, \texttt{sim/} and \texttt{walkthrough/}. Every statistical definition, gate and threshold in those files was specified by the authors; the model's contribution was implementation, review of implementation, and the identification of two defects that the authors then confirmed by inspection and by rerunning the affected analyses.
\item
  \textbf{Figure code.} All figures were drawn by R scripts written with model assistance (\texttt{figures/src/*.R}). Every number plotted was read from a released result table, and the correspondence between figure and table was checked by the authors panel by panel.
\item
  \textbf{Manuscript language and structure}: rewriting for length and readability, reorganisation of sections, and consistency checking of cross-references and reference numbering.
\end{itemize}

No text, number, figure or citation was accepted without author verification against the underlying result tables. The model did not select cohorts, choose thresholds, decide which analyses to report, or interpret results. Prompts were not used to generate data of any kind. The authors take full responsibility for the content of this publication.

\section{10. Data and code availability}\label{data-and-code-availability}

All cohorts are public. The lupus datasets are at GEO accessions GSE135779, GSE174188 and GSE285773; the COVID-19, influenza and CMV cohorts were obtained through the CELLxGENE Discover curation API and their dataset identifiers are in \texttt{cohorts/registry.yaml}. The MIT-licensed repository is at https://github.com/LightChainr/rheumlens with a permanent Zenodo record at https://doi.org/10.5281/zenodo.20813922. The versioned research object accompanying this manuscript contains the cohort registry with API-verified donor counts, the donor-level interface tables, every result table underlying Figures 2 to 8, all five per-seed outputs rather than a summary, the 19,600-cohort simulation output, analysis scripts, locked environments, \texttt{REPRODUCE.md} and \texttt{SHA256SUMS}. Public raw single-cell matrices remain at their original accessions.

\section{11. Declarations}\label{declarations}

\textbf{Author contributions.} Conceptualization, H.Y. and D.L.; methodology, H.Y.; software, H.Y.; validation, H.Y. and D.Y.; formal analysis, H.Y.; data curation, H.Y.; writing - original draft preparation, H.Y.; writing - review and editing, D.Y. and D.L.; visualization, H.Y.; supervision, D.L.; project administration, D.L. All authors have read and agreed to the submitted version of the manuscript.

\textbf{Funding.} This research received no external funding.

\textbf{Competing interests.} The authors declare no competing interests.

\textbf{Ethics approval.} Not applicable. This study analysed only publicly available, de-identified single-cell datasets released by their original investigators under the consent and approvals obtained at each contributing site. No new human or animal data were collected.

\textbf{Informed consent.} Not applicable, for the reason given above.

\textbf{Use of AI tools.} See Section 9.18, which states what was used, where, and how the output was checked.

\textbf{Acknowledgments.} The authors thank the investigators of GSE135779, GSE174188 and GSE285773, and the contributors to the COMBAT, Stephenson, Ren and CMV cohorts distributed through CELLxGENE Discover, for releasing donor-level data that made this re-analysis possible.

